\documentclass[12pt]{article}
\usepackage{lmodern}
\usepackage{xcolor}
\usepackage{amsmath, amsthm, amssymb}
\usepackage[margin=1in]{geometry}
\input{_macros}

\newtheorem{theorem}{Theorem}
\newtheorem{lemma}{Lemma}
\newtheorem{assumption}{Assumption}
\newtheorem{definition}{Definition}

\begin{document}

\section*{Section 3: Scientific Force Multiplication (C3)}

\subsection*{Assumption Used}

\begin{assumption}[Documented Scientific Breakthroughs (A5)]\label{A5}
  The following AI-driven scientific results are documented facts:
  \begin{enumerate}
    \item \textbf{AlphaFold2} (DeepMind, 2021): Solved the protein-structure
      prediction problem to near-experimental accuracy.  A5 explicitly states
      this problem had been \emph{open for 50 years} despite active human research.
      The resulting database covers $>$200M protein structures.
    \item \textbf{AI drug discovery} (Insilico Medicine, 2023): An AI-designed
      drug candidate reached IND approval in \textbf{18 months}.  The
      historical human-only average is \textbf{4--6 years} (48--72 months).
    \item \textbf{GNoME} (DeepMind, 2023): Discovered \textbf{2.2 million}
      novel stable crystal structures, surpassing the total of all prior
      human experimental discovery (approximately 20{,}000 structures) by a
      stated ratio of $>100\times$.  The structures are described as novel
      (not previously in the human record).
    \item \textbf{Lean4 AI proof assistance} (2023--2024): AI-assisted proof
      search produced \emph{novel proof strategies} for \emph{open}
      combinatorics problems; such tasks previously required years of
      specialist effort.
  \end{enumerate}
\end{assumption}

\subsection*{Definitions}

\begin{definition}[Prior Human Record]\label{def:record}
  The \emph{prior human record} in a domain is the collection of documented
  scientific results produced by human researchers before the AI system
  in question was deployed, as described in A5.
\end{definition}

\begin{definition}[Beyond-Record Contribution]\label{def:beyond}
  An AI result $R$ is a \emph{beyond-record contribution} in domain $D$
  if $R$ is explicitly described in A5 as either:
  (a)~a solution to a problem stated in A5 to have been \emph{open}
  (unsolved) in the prior human record, or
  (b)~a collection of results whose scale or novelty is stated in A5
  to exceed the entire prior human record in $D$.
\end{definition}

\begin{definition}[Time-to-Discovery Speedup]\label{def:speedup}
  For a discovery task $T$, the \emph{time-to-discovery speedup} is the
  ratio $\tau_{\mathrm{human}} / \tau_{\mathrm{AI}}$, where
  $\tau_{\mathrm{human}}$ is the documented human-only time and
  $\tau_{\mathrm{AI}}$ is the documented AI-assisted time for the same task,
  both as stated in A5.
\end{definition}

\begin{definition}[Independent Scientific Problem Areas]\label{def:independent}
  Problem areas are \emph{independent} if each is a distinct recognised
  scientific discipline and solving one does not constitute a solution to
  any other.  The four areas named in A5 --- protein-structure prediction
  (structural biology), drug candidate design (pharmacology/medicine),
  stable crystal-structure discovery (materials science), and formal
  combinatorics proof search (mathematics) --- are stipulated as
  independent by this definition: each belongs to a different branch of
  science, and no prior-record solution in one area is a solution in
  another.
\end{definition}

\subsection*{Statement to Prove (C3)}

\begin{theorem}[Scientific Force Multiplication]\label{thm:C3}
  Under Assumption~\ref{A5} and Definitions~\ref{def:record}--\ref{def:independent},
  AI systems constitute a cross-domain scientific force multiplier:
  \begin{enumerate}
    \item[(i)] In at least two independent scientific domains, AI has made
      beyond-record contributions (Definition~\ref{def:beyond}).
    \item[(ii)] A5 documents at least one time-to-discovery speedup
      $\geq 2\times$, and at least one beyond-record contribution whose
      scale exceeds the prior human record by $>100\times$.
    \item[(iii)] The documented gains span at least four independent
      scientific problem areas (Definition~\ref{def:independent}).
  \end{enumerate}
\end{theorem}

\begin{proof}

\textbf{Step 1: Domain inventory.}
A5 names four case studies in four distinct problem areas:
protein-structure prediction ($\alpha$),
drug candidate design ($\delta$),
stable crystal-structure discovery ($\gamma$), and
formal combinatorics proof search ($\lambda$).
By Definition~\ref{def:independent}, these four areas are independent.
This immediately establishes claim~(iii): the gains span four independent
problem areas.

\medskip
\noindent\textbf{Step 2: Beyond-record contributions in two domains.}

\emph{Case $\alpha$ (AlphaFold2).}
A5 states that protein-structure prediction was an \emph{open problem for
50 years}.  AlphaFold2 solved it.  A5's own language characterises this as
an open problem in the prior human record.  Therefore case $\alpha$ satisfies
Definition~\ref{def:beyond}(a).  Domain: structural biology.

\emph{Case $\gamma$ (GNoME).}
A5 states GNoME produced 2.2 million novel stable crystal structures and
explicitly gives the ratio over the prior human experimental corpus as
$>100\times$.  This directly satisfies Definition~\ref{def:beyond}(b):
the scale of the result exceeds the entire prior human record by
$>100\times$, as stated in A5.  Domain: materials science.

Two independent domains ($\alpha$ and $\gamma$) exhibit beyond-record
contributions.  This establishes claim~(i).

\medskip
\noindent\textbf{Step 3: Quantified acceleration.}

\emph{Time-to-discovery speedup (case $\delta$).}
By Definition~\ref{def:speedup}:
\[
  \tau_{\mathrm{human}} \;\geq\; 4 \text{ years} = 48\text{ months},
  \qquad
  \tau_{\mathrm{AI}} = 18\text{ months (A5)}.
\]
Using the conservative (shortest) human baseline from A5:
\[
  \text{Speedup}_\delta
  \;=\; \frac{\tau_{\mathrm{human}}}{\tau_{\mathrm{AI}}}
  \;\geq\; \frac{48}{18}
  \;=\; \frac{8}{3}
  \;\approx 2.67
  \;>\; 2.
\]
Hence A5 documents a time-to-discovery speedup $>2\times$.

\emph{Scale multiplication $>100\times$ (case $\gamma$).}
A5 explicitly states the GNoME discovery-count ratio is $>100\times$.
We adopt this A5-stated ratio directly.

Both quantified thresholds in claim~(ii) are established.

\medskip
\noindent\textbf{Conclusion.}
Claims (i), (ii), and (iii) are derived exclusively from A5 and
Definitions~\ref{def:record}--\ref{def:independent}.  Together they
establish that AI is a cross-domain scientific force multiplier with at
least two beyond-record contributions, one confirmed $>2\times$
time speedup, one confirmed $>100\times$ scale multiplication, spanning
four independent scientific disciplines.
\end{proof}

\end{document}
